% Обзор литературы 0004.
% Сборка: xelatex review.tex && bibtex review && xelatex review.tex && xelatex review.tex
% PDF кладётся в out/, в репозиторий не коммитится.
\documentclass[11pt,a4paper]{article}

\usepackage{iftex}
\ifPDFTeX
  \usepackage[T2A]{fontenc}
  \usepackage[utf8]{inputenc}
  \usepackage[russian]{babel}
\else
  \usepackage{fontspec}
  \usepackage[russian]{babel}
  \IfFontExistsTF{Times New Roman}
    {\setmainfont{Times New Roman}}
    {\IfFontExistsTF{DejaVu Serif}
      {\setmainfont{DejaVu Serif}}
      {\IfFontExistsTF{Liberation Serif}
        {\setmainfont{Liberation Serif}}
        {}}}
\fi

\usepackage[margin=2.4cm]{geometry}
\usepackage{amsmath}
\usepackage{booktabs}
\usepackage{longtable}
\usepackage{array}
\usepackage{enumitem}
\usepackage{hyperref}
\usepackage{microtype}

\hypersetup{
  colorlinks=true,
  linkcolor=black,
  citecolor=black,
  urlcolor=blue
}

\setlist[itemize]{leftmargin=1.4em,itemsep=0.2em,topsep=0.2em}
\setlist[enumerate]{leftmargin=1.6em,itemsep=0.2em,topsep=0.2em}

\title{Сборка датасета задач разной сложности,\\
включая нерешаемые задачи: обзор литературы}
\author{Обзор 0004}
\date{8 ноября 2026}

\begin{document}
\maketitle

\begin{abstract}
\noindent
Обзор отвечает на вопрос, как собрать набор задач с управляемым спектром
трудности, включая заведомо нерешаемые задачи, для оценки и обучения
бюджетного кодового агента. Включено 48 работ: стратификация готовых
бенчмарков, построение трудных наборов с нуля, генерация нерешаемых и
неотвечаемых задач, синтез задач с заданной трудностью, отбор данных по
трудности и аудит качества наборов. Вывод: готового рецепта нет, но есть
устойчивая схема. Трудность надо измерять на целевой паре «задача--агент», а
не брать из условия; нерешаемость — отдельная ось, которую нельзя выводить
из провала модели; у любой метки «нерешаемо» есть заметная доля ошибок, и
её нужно проверять независимо. Управляемая генерация даёт спектр трудности,
но платит за это падением выполнимости. Оптимальная пропорция трудных задач
зависит от бюджета данных, а покрытие распределения может быть важнее
средней трудности.
\end{abstract}

\section{Вопрос и зачем}

\subsection{Вопрос}

Как собрать набор задач с управляемым спектром сложности, включая заведомо
нерешаемые задачи, пригодный для оценки и обучения бюджетного кодового
агента?

Вопрос распадается на четыре независимых решения:

\begin{enumerate}
  \item \textbf{Откуда брать задачи} — существующий бенчмарк, реальные
        issue, синтетика или экспертная авторская работа.
  \item \textbf{Чем измерять трудность} — наблюдаемая доля успеха, латентная
        IRT-способность, внешняя метка платформы или экспертная оценка.
  \item \textbf{Как получать нерешаемые задачи} — писать вручную, портить
        условие, конфликтовать условие с тестами или брать задачи за
        границей возможностей модели.
  \item \textbf{Как проверить, что набор действительно имеет нужный спектр и
        что метка «нерешаемо» верна.}
\end{enumerate}

PICO:

\begin{itemize}
  \item \textbf{Population} — текстовые и кодовые задачи: Programming Online
        Judge, issue-resolution, вопросы на рассуждение, агентские задания.
  \item \textbf{Intervention} — метод построения набора: стратификация,
        генерация с заданной трудностью, генерация нерешаемых задач,
        проверка выполнимости.
  \item \textbf{Comparison} — случайная выборка из готового бенчмарка;
        набор одной трудности; набор без нерешаемых задач; набор без
        проверки выполнимости.
  \item \textbf{Outcome} — разброс трудности и разрешающая способность,
        точность метки «нерешаемо», устойчивость к контаминации, стоимость
        и воспроизводимость сборки.
\end{itemize}

\subsection{Зачем}

Харнессу нужен тренировочный набор: около 100 задач, 15--25\,\% из них
невыполнимы, баллы не связаны со сложностью, часть задач очень простая,
часть — очень трудная. Сейчас такого набора нет: примеры организаторов
содержат только выполнимые задачи, а финальный набор скрыт.

Ответ на вопрос определяет:
\begin{itemize}
  \item какие источники задач использовать и в какой пропорции;
  \item как назначать и проверять трудность;
  \item как строить нерешаемые задачи и сколько доверять этой метке;
  \item нужна ли спецификация изменения для сборки данных в \texttt{data/}
        и отдельный эксперимент по проверке набора.
\end{itemize}

Обзор не проектирует набор. Он фиксирует, какие процедуры уже проверены,
какие метрики использовались и где остаётся пробел.

\section{Метод}

\subsection{Базы и запросы}

Поиск выполнен 8 ноября 2026 года в двух базах: arXiv API (поле
\texttt{abs:}) и Crossref (поле \texttt{query.bibliographic}). Просматривались
верхние 20 записей на запрос, сортировка по релевантности. Язык —
английский, годы — 2015--2026.

Semantic Scholar вернул HTTP 403, OpenAlex исчерпал дневной лимит без
ключа. Поэтому вторая база — Crossref. У Crossref нет строгого совпадения
фраз: запрос без кавычек возвращает всю индексную выдачу, поэтому
\texttt{total-results} в миллионы записей не отражает число релевантных
работ. Для отбора считались только верхние 20 записей на запрос.

Точные запросы, даты, числа найденных и просмотренных записей приведены в
\texttt{protocol.md}. Выполнено 13 запросов к arXiv и 5 к Crossref.

\subsection{Критерии отбора}

Включались работы, которые:
\begin{itemize}
  \item описывают построение, отбор или разметку набора задач, а не только
        обучение модели на готовом наборе;
  \item дают количественную оценку свойства набора: разброс трудности,
        разрешающую способность, долю и качество нерешаемых задач,
        устойчивость к контаминации или стоимость сборки;
  \item относятся к тексту, коду, рассуждению или агентским задачам;
  \item доступны по ссылке (arXiv, DOI, страница издателя).
\end{itemize}

Исключались работы, где трудность задаётся только субъективно и без
проверки, где нерешаемость понимается как «модель не справилась», где нет
доступа к тексту или аннотации и где метод не переносится из узкого домена.

\subsection{Отбор}

Поиск вернул 3\,703 записи. Выдача ранжирована, поэтому просматривались
верхние 20 записей на запрос — всего 360. После удаления 16 дубликатов
осталось 344; на скрининге исключено 274 записи, не относящиеся к
построению или разметке набора. Для 70 работ прочитана аннотация или полный
текст, 22 исключены с причиной. В обзор включено 48 работ. Остальные 3\,343
записи за пределами верхних 20 не просматривались — это ограничение
обзора.

Схема и причины исключения — в \texttt{sources.md}.

\subsection{Извлечение и оценка качества}

Для каждой работы выписаны: что строится, источник задач, способ задания
трудности, способ получения нерешаемых задач, метрика и результат,
ограничение. Надёжность оценивалась по рецензируемости, открытости данных и
кода, числу задач, наличию процедуры разметки и согласия разметчиков,
отложенной проверке и воспроизводимости.

Мета-анализ не проводился: наборы, домены и метрики несопоставимы.

\section{Что известно}

\subsection{Трудность — свойство пары «задача--агент», а не текста}

Первое, что меняет постановку: численная трудность задачи не является
свойством только её условия.

Easy2Hard-Bench~\cite{easy2hard2024} собирает шесть наборов из математики,
кода, шахмат и рассуждений и назначает каждому заданию единую числовую
трудность через IRT и Glicko-2 по ответам людей или моделей. Шкала
получается согласованной внутри набора, но она относительна популяции
решателей: задача, трудная для одной модели, может быть лёгкой для другой.
TaskEval, разобранный в обзоре~0001, показывает то же для кода: трудность
зависит от формулировки промпта. Отсюда практическое следствие: метки
трудности, снятые с чужой модели или с рейтинга платформы, задают только
порядок, а не абсолютную вероятность успеха нашего агента.

Это не мешает использовать внешние метки для \emph{стратификации}: нам
нужен спектр, а не точная вероятность. Но калибровать границы «лёгкое /
среднее / трудное» надо на собственных прогонах.

\subsection{Три операциональные шкалы трудности}

Литература использует три разные шкалы, и они не взаимозаменяемы.

\textbf{Наблюдаемая доля успеха.} Самый простой вариант: прогнать модели и
разбить задачи по pass rate. Так устроены APPS~\cite{apps2021} (три уровня:
introductory, interview, competition), TACO~\cite{taco2023} (метки
сложности платформы и темы алгоритмов), LiveCodeBench~\cite{livecodebench2024}
(свежие контесты) и ZeroBench~\cite{zerobench2025} (adversarial filtering до
нуля успехов у фронтира). Шкала честная, но требует прогонов и устаревает
вместе с моделями.

\textbf{Латентная IRT-шкала.} Трудность и дискриминация оцениваются по
матрице ответов многих моделей. tinyBenchmarks~\cite{tinybenchmarks2024}
показывает, что 100 примеров достаточно, чтобы воспроизвести оценку MMLU на
14\,000. ATLAS~\cite{atlas2025} выбирает пункты по информации Фишера и
сокращает число пунктов до 90\,\% при той же точности: 41 пункт на HellaSwag
(5\,600 пунктов) даёт MAE $0{,}157$. IRT даёт не только трудность, но и
разрешающую способность — способность пункта различать близкие модели.

\textbf{Экспертная и структурная шкала.} GPQA~\cite{gpqa2023} опирается на
разрыв между экспертами (65\,\% правильных, 74\,\% без явных ошибок) и
неэкспертами (34\,\% после 30 минут поиска). HLE~\cite{hle2025} требует
однозначного проверяемого ответа, который нельзя быстро найти в интернете.
FrontierMath~\cite{frontiermath2024} — задачи, на которые у математика уходят
часы или дни. Здесь трудность задаётся не прогоном, а квалификацией
автора и процедурой проверки.

Для нашего набора разумна комбинация: внешняя метка или источник — для
грубого порядка, IRT по собственным прогонам — для калибровки и отбора,
экспертная проверка — для верхнего хвоста.

\subsection{Стратификация готовых наборов}

Самый дешёвый способ получить спектр — взять набор, где трудность уже
размечена, и выбрать задачи в нужных пропорциях.

Easy2Hard-Bench~\cite{easy2hard2024} — ближайшая к нам работа: единая
числовая шкала поверх шести разнородных наборов и намеренно повышенная доля
трудных задач. TACO~\cite{taco2023} даёт 25\,433 обучающих и 1\,000 тестовых
задач с темами, алгоритмами и уровнями сложности. APPS~\cite{apps2021} —
10\,000 задач с тремя уровнями и тестами. LiveCodeBench~\cite{livecodebench2024}
даёт 400 задач за окно май 2023 — май 2024 и снижает риск контаминации.

Ограничение этого пути очевидно: набор получается настолько трудным,
насколько его задумали авторы. Если нужна доля «очень простых» задач, её
придётся добирать из другого источника, а не вырезать из трудного набора.

\subsection{Сборка трудных наборов с нуля}

Когда готового набора нужной трудности нет, литература даёт несколько
проверенных процедур.

\textbf{Экспертная авторская работа с многоступенчатой проверкой.}
GPQA~\cite{gpqa2023} сначала пишут эксперты, затем проверяют неэксперты с
доступом в интернет: вопрос остаётся в наборе, только если неэксперт
ошибается, а эксперт отвечает правильно. HLE~\cite{hle2025} добавляет
требование однозначного решения и невозможности быстрого поиска.
FrontierMath~\cite{frontiermath2024} использует новые неопубликованные
задачи и автоматическую проверку ответа. Общий шаблон: \emph{критерии
качества формулируются до сбора}, а не после.

\textbf{Adversarial filtering.} ZeroBench~\cite{zerobench2025} отбирает
вопросы, на которых фронтирные модели дают 0\,\% при выпуске; через год
лучший результат — 6\,\% pass$^5$ и 19\,\% pass@5 на 46 моделях. Это
показывает и силу, и слабость подхода: набор действительно трудный, но
планка быстро сдвигается.

\textbf{Свежесть и контаминация.} LiveBench~\cite{livebench2024} обновляет
вопросы ежемесячно, использует объективную проверку и держит топовые модели
ниже 70\,\%. LiveCodeBench~\cite{livecodebench2024} собирает задачи после
даты обучения моделей. FrontierMath~\cite{frontiermath2024} держит задачи
неопубликованными. Это не про трудность напрямую, но без этого измеренный
спектр отражает память модели, а не способность решать.

\textbf{Проверяемость ответа.} BrowseComp~\cite{browsecomp2025} требует
короткий ответ, который легко сверить, хотя найти его трудно. HLE и GPQA
устроены так же. Для нашего набора это критично: если ответ нельзя
проверить автоматически, задача не годится для оценки.

\textbf{Реальные задачи как источник трудности.} SWE-bench~\cite{swebench2023}
берёт 2\,294 issue из 12 Python-репозиториев; сложность возникает из
необходимости менять несколько функций и файлов. BigCodeBench~\cite{bigcodebench2024}
строит 1\,140 задач на композицию вызовов функций с 5,6 тестами и покрытием
ветвей 99\,\%. Это ближе всего к нашему домену: трудность не выдумана, а
следует из структуры задачи.

\subsection{Нерешаемые задачи: сначала таксономия}

Слово «нерешаемая» покрывает разные явления, и методы их создания разные.
По включённым работам выделяются шесть классов.

\begin{enumerate}
  \item \textbf{Неотвечаемый вопрос при данном контексте.} SQuAD
        2.0~\cite{squad2_2018} добавляет 50\,000 adversarially написанных
        вопросов, ответа на которые в тексте нет. Неотвечаемость зависит от
        контекста, а не от задачи: тот же вопрос может быть отвечаем при
        другом документе.
  \item \textbf{Ложная предпосылка.} CREPE~\cite{crepe2022} собирает вопросы
        с форумов, где 25\,\% содержат ложную предпосылку; модель должна не
        отвечать, а исправлять предпосылку.
  \item \textbf{Недостающие условия.} В математике это задачи, где условие
        неполно и ответ не определён; Liu и соавторы~\cite{answering_unanswerable2025}
        показывают, что reasoning-модели видят дефект, но всё равно
        отвечают. ReliableMath~\cite{reliablemath2026} строит парные
        решаемые и нерешаемые задачи и находит, что признак
        «решаемости» переносится между языками, хотя верность ему падает.
  \item \textbf{Невозможность по конструкции.} The Impossible
        Test~\cite{impossible_test2024} — 675 graduate-level вопросов с
        намеренно неизвестным ответом. BSBench~\cite{bsbench2025} берёт
        существующие вопросы и модифицирует их так, чтобы разумного ответа
        не существовало. ZeroBench~\cite{zerobench2025} — «невозможный»
        визуальный набор, трудный для текущего фронтира.
  \item \textbf{Конфликт спецификации и проверки.} ImpossibleBench~\cite{impossiblebench2025}
        создаёт варианты LiveCodeBench и SWE-bench, где условие и unit-тесты
        прямо противоречат друг другу. Любое прохождение тестов означает
        обход спецификации; метрика — cheating rate.
  \item \textbf{За границей возможностей конкретной модели.} Unlocking the
        Unsolvable~\cite{unlocking_unsolvable2026} называет нерешаемыми
        задачи, на которых все rollout'ы модели проваливаются. Это не
        свойство задачи, а свойство пары «задача--модель»; 128 таких задач
        заменяют 2\,000 обычных при обучении (около 16$\times$ экономии).
\end{enumerate}

Есть и седьмой, пограничный класс: задачи, не решённые человечеством.
UQ~\cite{uq2025} собирает 500 нерешённых вопросов со Stack Exchange и
проверяет решения сообществом. Здесь «нерешаемость» неизвестна: задача
может быть решена завтра.

Для нашего набора важны первые пять классов. Шестой — это не нерешаемость
задачи, а сигнал обучения; седьмой — слишком дорог и не даёт стабильной
метки.

\subsection{Как получают нерешаемые задачи}

\textbf{Ручная авторская работа.} SQuAD 2.0~\cite{squad2_2018} и The
Impossible Test~\cite{impossible_test2024} пишутся людьми. Качество высокое,
цена — тоже; The Impossible Test содержит всего 675 задач.

\textbf{Автоматическая перестановка контекста.} AGent~\cite{agent_unanswerable2023}
берёт отвечаемый вопрос и подставляет контекст, где нужной информации нет.
Так получаются наборы из SQuAD и HotpotQA с низкой долей ошибок, а
дообучение на них сравнимо с SQuAD 2.0.

\textbf{Дешёвые замены.} Gautam и соавторы~\cite{lightweight_unanswerable2023}
меняют антоним или сущность в отвечаемом вопросе. Метод без обучения даёт
+1,6 F1 на SQuAD 2.0 и +9,3 F1 на TydiQA с BERT-large. Это сильный baseline:
сложная генерация не всегда оправдана.

\textbf{LLM-генерация с проверкой.} CoDiQ~\cite{codiq2026} и MathSmith~\cite{mathsmith2025}
генерируют задачи и проверяют согласованность ответа. Важно, что
генератор оптимизирует не только трудность, но и решаемость: CoDiQ получает
44\,000 задач с более чем 82\,\% решаемости, труднее LiveCodeBench и AIME.

\textbf{Модификация тестов.} ImpossibleBench~\cite{impossiblebench2025}
не пишет задачи, а меняет проверку. Это самый дешёвый способ получить
«нерешаемую» версию готового набора, но он измеряет склонность к обходу,
а не трудность.

\textbf{Комбинирование.} AbstentionBench~\cite{abstentionbench2025}
объединяет 20 наборов и пять причин отказа: неизвестный ответ,
недоопределённость, ложная предпосылка, субъективная интерпретация и устаревшие
данные. Это удобно для
оценки, но плохо для обучения: причина отказа смешана, и модель не учится
различать классы.

\subsection{Проверка выполнимости: где ломаются метки}

Главный результат группы — метка «нерешаемо» сама по себе ненадёжна.

SWE-bench+~\cite{swebenchplus2024} вручную проверяет успешные патчи и
находит: 32{,}67\,\% решений прямо утекают из текста issue, а 31{,}08\,\%
проходят слабые тесты. После фильтрации resolve rate падает с 12{,}47\,\%
до 3{,}97\,\%. Более 94\,\% задач созданы до даты отсечения знаний модели.
То есть часть «решённых» задач была нечестно решена, а часть «провалов» —
следствие плохих тестов.

Who Guards the Benchmarks?~\cite{whoguards2026} автоматически аудитирует
исполняемые агентские наборы и находит в ScienceAgentBench 12
подтверждённых авторами проблем, включая фатальные ошибки, делающие задачи
нерешаемыми. На BIXBench Verified-50 совпадение с экспертами — 83{,}3\,\%,
стоимость аудита 50 задач — меньше \$15.

Are We Done with MMLU?~\cite{mmlu_redux2024} переаннотирует 5\,700 пунктов
и оценивает долю ошибок в 6{,}49\,\% (в вирусологии — 57\,\%). Auditing LLM
Benchmarks with IRT~\cite{audit_irt2026} находит вероятные ошибочные метки с
точностью 95\,\% в топ-200 на семи наборах по ответам 114 моделей.

Отсюда три обязательных проверки для наших нерешаемых задач:
\begin{enumerate}
  \item \textbf{Независимый решатель.} Сильная модель с доступом к тестам и
        среде не должна решать задачу; иначе она не нерешаемая.
  \item \textbf{Аудит тестов.} Тесты должны покрывать спецификацию; слабые
        тесты превращают решаемую задачу в «нерешаемую» и наоборот.
  \item \textbf{Проверка спецификации.} Условие не должно содержать ответ
        или противоречие, которое делает задачу бессмысленной. Для этого
        годится LLM-аудитор, но с выборочной человеческой проверкой.
\end{enumerate}

\subsection{Генерация с управляемой трудностью}

Если готовых задач нужной трудности нет, её можно синтезировать. Литература
даёт четыре рычага.

\textbf{Явная метка и переписывание.} DQG~\cite{dqg2018} генерирует вопросы
с заданной меткой трудности; Guiding the Growth~\cite{guiding_growth2021}
определяет трудность как число шагов вывода и повышает её пошаговым
переписыванием. Обе работы показывают, что трудностью можно управлять, но
метка остаётся прокси.

\textbf{Test-time scaling.} CoDiQ~\cite{codiq2026} увеличивает бюджет
рассуждения генератора и обнаруживает важный компромисс: более длинное
рассуждение повышает трудность, но снижает решаемость. Это прямое
предупреждение: оптимизировать трудность без ограничения на выполнимость
нельзя.

\textbf{Декомпозиция и стратегии.} MathMixup~\cite{mathmixup2026} разлагает
задачу на компоненты и собирает градацию трудности; Qwen2.5-7B после
обучения на таком наборе даёт 52{,}6\,\% в среднем на семи математических
бенчмарках. MathSmith~\cite{mathsmith2025} строит задачи с нуля из
концептов PlanetMath и применяет девять стратегий усложнения и RL по длине
рассуждения. SPARQ~\cite{sparq2025} использует quality-diversity и
solve-rate как прокси трудности: из 7,5 тысяч затравок получается 20 млн пар,
а фильтрация по трудности даёт до +24\,\% относительного прироста.

\textbf{Предсказание трудности без прогонов.} R2DE~\cite{r2de2020} оценивает
IRT-параметры нового вопроса по тексту и вариантам ответа, без предтеста.
Ballon и соавторы~\cite{llmcompare2025} идут дальше: попарные сравнения
задач языковой моделью и модель Брэдли--Терри дают шкалу без эталонных
ответов, с корреляцией Пирсона не ниже 0{,}80 на 1\,876 парах и деградацией
меньше 6\,\% при 10\,\% шума. Это единственный найденный способ назначить
трудность задачам, которые пока не решает никто.

Общее ограничение группы: синтетика оптимизирует правдоподобие и трудность,
но не гарантирует, что задача вообще имеет решение. Поэтому генератор
должен сопровождаться верификатором, а доля решаемых задач — измеряться, а
не предполагаться.

\subsection{Сколько трудности нужно и в какой пропорции}

Для оценки и для обучения ответы разные.

\textbf{Для оценки} нужны пункты с высокой разрешающей способностью вокруг
текущего уровня агента. Это аргумент IRT: ATLAS~\cite{atlas2025} и
tinyBenchmarks~\cite{tinybenchmarks2024} отбирают именно такие пункты и
сокращают набор без потери точности ранжирования. Но для оценки
\emph{предела} нужен и верхний хвост: ZeroBench~\cite{zerobench2025},
FrontierMath~\cite{frontiermath2024} и HLE~\cite{hle2025} специально состоят
из задач, которые текущие модели почти не решают.

\textbf{Для обучения} оптимальная трудность не универсальна.
Liu и соавторы~\cite{difficulty_tradeoff2026} показывают, что при
фиксированном бюджете данных существует оптимальная трудность, и она
сдвигается к более трудным примерам по мере роста бюджета. Gao и
соавторы~\cite{selective_dpo2025} находят, что слишком трудные примеры
вредят alignment: фильтрация даёт +9--16\,\% win rate на AlpacaEval 2.
Pandia и соавторы~\cite{coverage2026} идут ещё дальше: для fine-tuning LLM
случайный отбор часто лучше difficulty-based, потому что важнее покрытие
распределения, а не средняя трудность. ZPD Detector~\cite{zpd2026}
пытается совместить оба взгляда: IRT-трудность сопоставляется с текущей
способностью модели, и отбираются задачи в зоне ближайшего развития.

Для нашего случая это означает: нельзя требовать «побольше трудных задач».
Нужен явно заданный спектр и проверка, что каждая полоса действительно
покрыта. Для тренировочного набора разумно держать большое число простых и
средних задач, небольшую полосу трудных и отдельную полосу нерешаемых; для
оценки — наоборот, больше задач вокруг границы способности.

\subsection{Аудит набора и контаминация}

Качество набора — не только трудность. BetterBench~\cite{betterbench2024}
формулирует 46 практик и оценивает 24 бенчмарка; у многих нет отчёта о
статистической значимости и воспроизводимости. Who Guards~\cite{whoguards2026}
показывает, что ошибки инфраструктуры выглядят как провалы агента.
Are We Done with MMLU~\cite{mmlu_redux2024} и IRT-аудит~\cite{audit_irt2026}
дают методы поиска ошибочных меток.

Контаминация решается тремя приёмами: свежесть (LiveBench, LiveCodeBench),
неопубликованность (FrontierMath) и временные окна. Для нашего набора,
который собирается из открытых задач, полностью избежать утечки нельзя;
можно её измерить и ограничить, проверив, как модель отвечает на
переформулировки и на задачи после даты отсечения.

\section{Сравнение}

Таблица~\ref{tab:sources} сводит включённые работы по тому, что они строят,
как задают трудность, как получают нерешаемость и какое у них ограничение.
Строки сгруппированы по назначению: стратификация, трудные наборы,
нерешаемые задачи, управляемая генерация, распределение для обучения.

{\footnotesize
\setlength{\tabcolsep}{3pt}
\renewcommand{\arraystretch}{1.12}
\begin{longtable}{@{}p{0.65cm}p{2.6cm}p{2.6cm}p{2.9cm}p{2.9cm}p{2.7cm}@{}}
\caption{Включённые работы.}\label{tab:sources}\\
\toprule
ID & Работа & Что строит & Трудность & Нерешаемость & Результат и ограничение \\
\midrule
\endfirsthead
\toprule
ID & Работа & Что строит & Трудность & Нерешаемость & Результат и ограничение \\
\midrule
\endhead
\bottomrule
\endfoot

S01 & Ding et al. \cite{easy2hard2024} & метки для 6 наборов & IRT, Glicko-2 & нет & единая шкала; нужна матрица ответов \\
S02 & Li et al. \cite{taco2023} & задачи по алгоритмам & метка платформы & нет & 25\,433 задачи; метка — прокси \\
S03 & Hendrycks et al. \cite{apps2021} & генерация кода & 3 уровня источника & нет & 10\,000 задач; уровни не калиброваны \\
S04 & Jain et al. \cite{livecodebench2024} & свежий набор & контест и свежесть & нет & 400 задач; метки трудности нет \\
S05 & Polo et al. \cite{tinybenchmarks2024} & сжатые бенчмарки & IRT & нет & 100 примеров вместо 14\,000 \\
S06 & Li et al. \cite{atlas2025} & адаптивный тест & IRT, информация Фишера & нет & экономия до 90\,\% пунктов; нужна калибровка \\
S07 & Gema et al. \cite{mmlu_redux2024} & переразметка & ручная проверка & ошибки меток & 6{,}49\,\% ошибок; один бенчмарк \\
S08 & Reuel et al. \cite{betterbench2024} & аудит & нет & нет & 46 практик, 24 бенчмарка \\
S09 & Land, Bikel \cite{audit_irt2026} & детектор меток & IRT-остатки & неоднозначные пункты & 95\,\% точности в топ-200 \\
S10 & Tu et al. \cite{whoguards2026} & аудит агентских наборов & нет & фатальные ошибки & 12 проблем; < \$15 на 50 задач \\
S11 & Jimenez et al. \cite{swebench2023} & реальные issue & структура задачи & нет & 2\,294 задачи; 1{,}96\,\% у Claude 2 \\
S12 & Aleithan et al. \cite{swebenchplus2024} & аудит SWE-bench & нет & слабые тесты & 31{,}08\,\%; resolve 12{,}47 $\to$ 3{,}97\,\% \\
S13 & Yang et al. \cite{swesmith2025} & синтетические issue & нет & нет & 50\,000 задач; 40{,}2\,\% на Verified \\
S14 & Zhuo et al. \cite{bigcodebench2024} & сложные задачи & композиция вызовов & нет & 99\,\% покрытия; человек 97\,\% \\
S15 & Rein et al. \cite{gpqa2023} & экспертный Q\&A & разрыв эксперт/неэксперт & нет & 448 вопросов; GPT-4 39\,\% \\
S16 & Phan et al. \cite{hle2025} & экспертный Q\&A & эксперт, запрет поиска & нет & 2\,500 вопросов; низкая калибровка \\
S17 & Glazer et al. \cite{frontiermath2024} & математика & часы-дни эксперта & нет & SOTA < 2\,\%; закрытый набор \\
S18 & Wei et al. \cite{browsecomp2025} & браузерные задачи & трудный поиск & нет & 1\,266 вопросов; ответ проверяем \\
S19 & Mialon et al. \cite{gaia2023} & реальные задачи & 3 уровня & нет & 466 вопросов; человек 92\,\% \\
S20 & Roberts et al. \cite{zerobench2025} & «невозможный» набор & adversarial filtering & для текущего фронтира & 0\,\% при выпуске; 6\,\% через год \\
S21 & White et al. \cite{livebench2024} & обновляемый набор & свежесть & нет & топ < 70\,\%; ежемесячно \\
S22 & Nie et al. \cite{uq2025} & нерешённые вопросы & неизвестна & не решены людьми & 500 вопросов; топ проходит 15\,\% \\
S23 & Rajpurkar et al. \cite{squad2_2018} & неотвечаемые вопросы & нет & нет ответа в контексте & 50\,000 вопросов; зависит от контекста \\
S24 & Tran et al. \cite{agent_unanswerable2023} & генерация неотвечаемых & нет & перестановка контекста & низкая доля ошибок \\
S25 & Gautam et al. \cite{lightweight_unanswerable2023} & аугментация & нет & замена антонима & +1,6 F1; +9,3 F1 на TydiQA \\
S26 & Yin et al. \cite{selfaware2023} & неотвечаемые вопросы & нет & 5 видов незнания & 20 LLM; разрыв с человеком \\
S27 & Yu et al. \cite{crepe2022} & ложные предпосылки & нет & 25\,\% вопросов & нужны исправления предпосылок \\
S28 & Kirichenko et al. \cite{abstentionbench2025} & сводный набор & нет & 5 причин отказа & reasoning ухудшает отказ на 24\,\% \\
S29 & Noever, McKee \cite{impossible_test2024} & невозможные вопросы & нет & ответ неизвестен & 675 задач; 62--68\,\% признания \\
S30 & Erziev \cite{bsbench2025} & модификация наборов & нет & нет разумного ответа & модели далеки от идеала \\
S31 & Zhong et al. \cite{impossiblebench2025} & конфликт условия и тестов & нет & обход спецификации & cheating rate \\
S32 & Miyai et al. \cite{upd2024} & детекция нерешаемых MCQ & нет & 3 типа & модели плохо удерживают ответ \\
S33 & Zoumpoulidi et al. \cite{reliablemath2026} & парные задачи & нет & пара решаемая/нерешаемая & признак переносится между языками \\
S34 & Liu et al. \cite{answering_unanswerable2025} & анализ отказа & нет & недостающие условия & мониторинг + вмешательство \\
S35 & Zhu, Han \cite{unlocking_unsolvable2026} & обучающий набор & трассы учителя & провал всех rollout & 128 задач $\approx$ 2\,000 \\
S36 & Gao et al. \cite{dqg2018} & генератор вопросов & явная метка & нет & соблюдает метку; метка субъективна \\
S37 & Cheng et al. \cite{guiding_growth2021} & усложнение & число шагов вывода & нет & трудность = длина цепочки \\
S38 & Peng et al. \cite{codiq2026} & трудные вопросы & test-time scaling & контроль решаемости & 44\,000 задач; > 82\,\% решаемы \\
S39 & Li et al. \cite{mathmixup2026} & синтез & декомпозиция & самопроверка & 52{,}6\,\% на 7 бенчмарках \\
S40 & Zhan et al. \cite{mathsmith2025} & синтез & 9 стратегий + RL & проверка ответа & лучше на AIME, OlympiadBench \\
S41 & Havrilla et al. \cite{sparq2025} & quality-diversity & solve-rate & нет & 20 млн пар; +24\,\% \\
S42 & Benedetto et al. \cite{r2de2020} & IRT-оценка & по тексту вопроса & нет & без предтеста; MCQ \\
S43 & Ballon et al. \cite{llmcompare2025} & шкала без эталона & попарные сравнения & нет & $r \geq 0{,}80$; устойчиво к шуму \\
S44 & Lalor, Yu \cite{ddaclae2020} & отбор данных & выученная трудность & нет & лучше эвристик на GLUE \\
S45 & Liu et al. \cite{difficulty_tradeoff2026} & анализ & оптимальная трудность & нет & растёт с бюджетом данных \\
S46 & Gao et al. \cite{selective_dpo2025} & фильтрация & трудность примера & нет & +9--16\,\% win rate \\
S47 & Pandia et al. \cite{coverage2026} & отбор данных & критика трудности & нет & покрытие важнее трудности \\
S48 & Yang et al. \cite{zpd2026} & отбор данных & IRT + способность & нет & зона ближайшего развития \\
\end{longtable}
}

\section{Пробел и вывод}

\subsection{Что не покрыто}

\begin{enumerate}
  \item \textbf{Единого рецепта смешанного набора нет.} Работы делают одно из
        трёх: стратифицируют готовый набор~\cite{easy2hard2024,taco2023},
        строят трудный набор с нуля~\cite{gpqa2023,hle2025,frontiermath2024}
        или генерируют нерешаемые задачи~\cite{squad2_2018,impossible_test2024,impossiblebench2025}.
        Ни одна не задаёт целевые доли «простое / среднее / трудное /
        нерешаемое» и не проверяет, что после сборки эти доли сохранились.
  \item \textbf{Нет метрики качества спектра.} Число задач и средняя
        трудность мало что говорят. Нужны разрешающая способность по
        полосам, покрытие диапазона способности и точность метки
        «нерешаемо». IRT даёт инструменты~\cite{atlas2025,tinybenchmarks2024},
        но стандарта для смешанного набора нет.
  \item \textbf{Метка «нерешаемо» плохо проверена на масштабе.} SQuAD
        2.0 и The Impossible Test — небольшие или узкие; SWE-bench+
        показывает, что и в крупных наборах 31\,\% проверок слабые. Работы,
        которые измеряют нерешаемость агентских задач, часто определяют её
        как провал модели~\cite{unlocking_unsolvable2026}, а это другая
        величина.
  \item \textbf{Генерация с управляемой трудностью почти не проверяет
        выполнимость.} CoDiQ~\cite{codiq2026} — исключение: он явно
        измеряет долю решаемых задач. Большинство методов сообщают только
        трудность и качество ответа.
  \item \textbf{Пропорция трудных задач для обучения зависит от бюджета и
        модели.} Общие рекомендации противоречат друг другу:
        \cite{difficulty_tradeoff2026,selective_dpo2025} за трудные примеры,
        \cite{coverage2026} против. Для нашего маленького набора (сотни
        задач) это надо мерить, а не выбирать по литературе.
  \item \textbf{Кодовые агентские наборы почти не размечены по трудности.}
        SWE-bench и LiveCodeBench дают сложность через источник, а не через
        калибровку. Easy2Hard-Bench покрывает код, но не issue-resolution.
\end{enumerate}

\subsection{Что из этого следует для сборки набора}

Литература не даёт готового набора, но даёт порядок действий.

\begin{enumerate}
  \item \textbf{Сначала определить целевой спектр, потом источники.} Для
        тренировочного набора задать доли явно: например, 15\,\% очень
        простых, 25\,\% простых, 25\,\% средних, 15\,\% трудных и
        20\,\% нерешаемых. Доли — гипотеза, которую надо проверить на
        собственных прогонах; в сумме они дают 100\,\%.
  \item \textbf{Брать задачи из четырёх источников.}
        \begin{itemize}
          \item готовые наборы с метками трудности: TACO, APPS,
                Easy2Hard-Bench;
          \item свежие задачи: LiveCodeBench, LiveBench;
          \item реальные issue: SWE-bench-подобные задачи с тестами;
          \item синтетика с управляемой трудностью: CoDiQ, MathMixup,
                SPARQ.
        \end{itemize}
        Смешивание источников даёт и разброс, и разные типы трудности.
  \item \textbf{Калибровать трудность на своём агенте.} Внешняя метка —
        только порядок. После пилота прогнать агента на подвыборке,
        оценить IRT-трудность и разрешающую способность, пересобрать
        полосы. Если прогонов мало, использовать попарные сравнения
        LLM~\cite{llmcompare2025} как дешёвую промежуточную шкалу.
  \item \textbf{Строить нерешаемые задачи четырьмя способами.}
        \begin{itemize}
          \item недостающие условия (математика, код без входных данных);
          \item ложная предпосылка (вопрос про несуществующую функцию или
                свойство);
          \item конфликт условия и тестов (ImpossibleBench-подобный);
          \item задача за границей возможностей модели (по результатам
                пилота).
        \end{itemize}
        Первые три дают стабильную метку, четвёртая — обучающий сигнал, а
        не свойство задачи. Для оценки нужны первые три.
  \item \textbf{Проверять каждую нерешаемую задачу тремя воротами.}
        \begin{enumerate}
          \item независимый сильный решатель с доступом к тестам не решает;
          \item аудит тестов: покрытие спецификации, отсутствие ответа в
                условии, отсутствие противоречий;
          \item выборочная человеческая или LLM-проверка формулировки.
        \end{enumerate}
        Без этого доля ошибок будет как в SWE-bench+ и MMLU: проценты, а не
        доли процента.
  \item \textbf{Измерять не только трудность, но и покрытие.} Для каждой
        полосы считать число задач, разрешающую способность и разброс
        стоимости в токенах. Если полоса пустая или все задачи в ней
        решаются одинаково, набор не различает агентов.
  \item \textbf{Ограничить контаминацию.} Предпочитать свежие и
        синтетические задачи, проверять переформулировки, держать часть
        набора закрытой. Полностью избежать утечки из открытых источников
        нельзя.
  \item \textbf{Проверить набор экспериментом.} Сначала пилот на 100
        задачах: измерить, как меняется выгода на токен при переходе от
        текущего порядка к стратифицированному набору. Спецификация
        изменения поведения нужна только если прирост устойчив.
\end{enumerate}

\subsection{Схема набора}

Таблица~\ref{tab:recipe} переводит выводы в конкретную схему для набора из
100 задач. Доли — гипотеза: их надо проверить на пилоте и сдвинуть, если
полоса не различает агента или стоит слишком дорого.

{\small
\begin{table}[h]
\centering
\caption{Схема тренировочного набора.}\label{tab:recipe}
\begin{tabular}{@{}p{2.6cm}p{1.5cm}p{3.2cm}p{3.4cm}p{3.6cm}@{}}
\toprule
Полоса & Доля & Источник & Как получить & Проверка \\
\midrule
Очень простые & 15 (10--20) & APPS introductory, TACO low & взять готовые метки & агент решает почти всегда \\
Простые & 25 (20--30) & TACO mid, APPS interview & готовые метки + IRT & успех выше среднего \\
Средние & 25 (20--30) & LiveCodeBench, простые SWE-bench & свежие задачи и реальные issue & успех около границы \\
Трудные & 15 (10--20) & TACO high, SWE-bench, CoDiQ, MathSmith & метки платформы + синтез & успех низкий, но не нулевой \\
Нерешаемые & 20 (15--25) & недостающие условия, ложные API, ImpossibleBench, задачи вне возможностей & 4 генератора & 3 ворот: решатель, аудит тестов, проверка формулировки \\
\bottomrule
\end{tabular}
\end{table}
}

Порядок сборки: сначала набрать полосы из готовых источников, затем
добавить синтетику и нерешаемые задачи, затем прогнать пилот и
перекалибровать границы по IRT. Если после проверки доля нерешаемых
проседает, добрать их и проверить заново: ошибка метки здесь дороже, чем
нехватка задач.

\subsection{Короткий ответ}

Собрать набор разной трудности можно, но не одной процедурой. Нужно
смешать четыре источника, задать целевые доли полос и откалибровать
трудность IRT по собственным прогонам, потому что трудность — свойство пары
«задача--агент». Нерешаемые задачи — отдельная ось: их надо строить
целенаправленно (недостающие условия, ложные предпосылки, конфликт с
тестами) и проверять независимым решателем и аудитом тестов, а не выводить
из провала модели. Управляемая генерация даёт нужный спектр, но снижает
выполнимость, поэтому долю решаемых задач надо измерять. Для обучения
оптимальная доля трудных задач зависит от бюджета данных, а покрытие
распределения может быть важнее средней трудности. Готового набора под
бюджетного кодового агента литература не даёт — его надо собрать и
проверить экспериментально.

\section{Ограничения обзора}

\begin{itemize}
  \item Две базы вместо четырёх: Semantic Scholar вернул HTTP 403, OpenAlex
        исчерпал дневной лимит. Crossref плохо ранжирует CS-препринты,
        поэтому почти все включённые работы найдены в arXiv. Возможны
        пропуски, особенно в индустриальных публикациях.
  \item Просматривались верхние 20 записей на запрос, а не вся выдача.
  \item Для большинства работ данные извлечены из аннотаций; полные тексты
        читались выборочно. Числа могут отличаться от итоговых таблиц
        статей.
  \item Препринты 2026 года не рецензированы и могут измениться.
  \item Метрики несопоставимы: число задач, доля нерешаемых, корреляция,
        разрешающая способность, стоимость сборки. Мета-анализ невозможен.
  \item Обзор не проверяет воспроизводимость сборки наборов и не оценивает
        публикационное смещение.
\end{itemize}

\bibliographystyle{plain}
\bibliography{refs}

\end{document}
